%% kapa: exam=1 round=33 year=2022 period=1 form=A subject=civil_law subject_ko=민법 seq=30 number=30
%% source: Q-Net 감정평가사 기출문제 (cst00309) artlSeq=5209997 file=raw/1008/5209997/2022년도 제33회 감정평가사 1차시험 1교시(민법,경제학,부동산학).pdf
%% conversion: kordoc 4.24.51 -> tools/gen1.py (auto); review=proofread-sonnet; flags=none
\begin{question}{30}
관습상의 법정지상권에 관한 설명으로 옳지 않은 것은? (다툼이 있으면 판례에 따름)
\begin{choices}
  \item 토지 또는 그 지상 건물의 소유권이 강제경매절차로 인하여 매수인에게 이전된 경우, 매수인의 매각대금 완납시를 기준으로 토지와 그 지상 건물이 동일인 소유에 속하였는지 여부를 판단하여야 한다.
  \item 관습상의 법정지상권이 성립하였으나 건물 소유자가 토지 소유자와 건물의 소유를 목적으로 하는 토지 임대차계약을 체결한 경우, 그 관습상의 법정지상권은 포기된 것으로 보아야 한다.
  \item 관습상의 법정지상권은 이를 취득할 당시의 토지소유자로부터 토지소유권을 취득한 제3자에게 등기없이 주장될 수 있다.
  \item 관습상의 법정지상권이 성립한 후에 건물이 증축된 경우, 그 법정지상권의 범위는 구 건물을 기준으로 그 유지·사용을 위하여 일반적으로 필요한 범위 내의 대지 부분에 한정된다.
  \item 관습상의 법정지상권 발생을 배제하는 특약의 존재에 관한 주장·증명책임은 그 특약의 존재를 주장하는 측에 있다.
\end{choices}
\end{question}
\answer{1}
